% =====================================================================
%  1. OBJETIVO
% =====================================================================
\section{Objetivo}
Implementar un programa de \matlab{} que combine la planificaci\'{o}n global (Dijkstra de Lab-MR 4) con la navegaci\'{o}n local reactiva por campos potenciales (Lab-MR 3) sobre el mapa \texttt{mapa2.pgm} y el grafo \texttt{mapa2.m}. El programa pide los nodos de inicio y destino, representa gr\'{a}ficamente el mapa con la trayectoria del robot e indica si se ha alcanzado el destino. Adicionalmente: realizar la trayectoria desde el nodo 1 hacia la zona superior derecha y atravesar la regi\'{o}n entre los nodos 24 y 32 (donde la navegaci\'{o}n b\'{a}sica falla por m\'{i}nimo local), proponer una mejora, y reemplazar Dijkstra por A* con una heur\'{i}stica eucl\'{i}dea consistente.


% =====================================================================
%  2. IMPLEMENTACIÓN
% =====================================================================
\section{Implementaci\'{o}n}

La soluci\'{o}n es un \'{u}nico script (\texttt{navegacion\_autonoma\_p7.m}) que carga el mapa con \texttt{binaryOccupancyMap}, ejecuta \texttt{mapa2.m} para obtener \texttt{nodos} y \texttt{costes}, planifica la ruta global y recorre los subobjetivos con un bucle local de hasta $1200$ iteraciones por tramo. Las decisiones no obvias son:

\paragraph{Matriz de costes eucl\'{i}dea para A*.} A* solo es \'{o}ptimo con heur\'{i}stica eucl\'{i}dea si los pesos del grafo coinciden con la m\'{e}trica eucl\'{i}dea. Por eso, antes de invocar \texttt{aestrella}, se construye una matriz $C^{\text{eucl}}$ que mantiene la topolog\'{i}a (adyacencia) del grafo original pero sustituye cada peso $C(i,j)$ por $\lVert p_i - p_j\rVert_2$. As\'{i} la heur\'{i}stica $H(i,j) = \lVert p_i - p_j\rVert_2$ es consistente.

\paragraph{Campo tangencial para mitigar m\'{i}nimos locales.} Adem\'{a}s de los campos atractivo y repulsivo de Lab-MR 3, se a\~{n}ade un campo tangencial $\mathbf{F}_\text{tan}$ orientado perpendicularmente al obst\'{a}culo m\'{a}s cercano, hacia el lado en el que se encuentra el subobjetivo (signo del producto vectorial entre la direcci\'{o}n al objetivo y la direcci\'{o}n al obst\'{a}culo). Su magnitud se modula por la cercan\'{i}a $\max(0,\, 1 - d/D)$. Este campo desempata las situaciones de equilibrio entre atracci\'{o}n y repulsi\'{o}n.

\paragraph{Detecci\'{o}n de estancamiento por progreso.} Se mantiene una ventana deslizante con las \'{u}ltimas $W = 40$ distancias al subobjetivo. Si el progreso acumulado no supera un umbral $\Delta_\text{min}$, se considera estancamiento. Este criterio es m\'{a}s robusto que medir velocidad nula porque el robot puede oscilar lateralmente a velocidad alta sin acercarse al objetivo.

\paragraph{Maniobra de escape.} Al detectar estancamiento, se prueban $K = 16$ direcciones equiespaciadas alrededor de la direcci\'{o}n al subobjetivo, se descartan las que conducen a colisi\'{o}n y se selecciona la que maximiza una puntuaci\'{o}n combinada de despeje y cercan\'{i}a al subobjetivo. Tras un n\'{u}mero m\'{a}ximo de estancamientos consecutivos, el tramo se aborta como m\'{i}nimo local.

\paragraph{Detecci\'{o}n de colisi\'{o}n por est\'{e}ncil de 9 puntos.} Antes de cada avance se comprueba colisi\'{o}n no en un \'{u}nico punto, sino en el centro y en 8 puntos circundantes a distancia \texttt{collision\_margin}. Esto modela un radio de seguridad y previene aproximaciones a esquinas que el campo continuo no detecta a tiempo.


% =====================================================================
%  3. PRUEBAS
% =====================================================================
\section{Pruebas}

\subsection{Caso base: nodo 1 a la parte superior derecha}
Con Dijkstra como planificador global y los campos potenciales b\'{a}sicos (sin campo tangencial ni maniobra de escape), el robot recorre la mayor\'{i}a de los tramos sin incidencias y alcanza el destino en la zona superior derecha del mapa. Activando la mejora se obtiene un comportamiento equivalente, sin estancamientos, con una traza pr\'{a}cticamente id\'{e}ntica.

\begin{figure}[H]
  \centering
  \begin{subfigure}[t]{0.48\linewidth}
    \centering
    \includegraphics[width=\linewidth]{nav_dijkstra_1_a_12_sin_mejora.jpg}
    \caption{Sin mejora.}
    \label{fig:nav_1_12_sin}
  \end{subfigure}\hfill
  \begin{subfigure}[t]{0.48\linewidth}
    \centering
    \includegraphics[width=\linewidth]{nav_dijkstra_1_a_12_con_mejora.jpg}
    \caption{Con mejora.}
    \label{fig:nav_1_12_con}
  \end{subfigure}
  \caption{Caso base nodo 1 $\rightarrow$ nodo 12 con planificaci\'{o}n global Dijkstra y navegaci\'{o}n local por campos potenciales sobre \texttt{mapa2.pgm}. Aristas del grafo en gris, ruta global en amarillo discontinuo, trayectoria del robot en rojo, origen verde y destino magenta.}
  \label{fig:nav_caso_base}
\end{figure}

\subsection{Zona conflictiva entre nodos 24 y 32}
Al intentar atravesar la regi\'{o}n entre los nodos 24 y 32, la versi\'{o}n b\'{a}sica del controlador queda atrapada de forma reproducible: el pasillo es estrecho y los obst\'{a}culos est\'{a}n dispuestos sim\'{e}tricamente respecto al eje del robot, con lo que la atracci\'{o}n al subobjetivo y la repulsi\'{o}n combinada se cancelan. Es la situaci\'{o}n de m\'{i}nimo local cl\'{a}sica que ya se identificaba en Lab-MR 3, agravada por la geometr\'{i}a del pasillo.

\textbf{Mejora propuesta:} a\~{n}adir el campo tangencial $\mathbf{F}_\text{tan}$ (que rompe la simetr\'{i}a empujando al robot hacia el lado del subobjetivo) y, como red de seguridad, una maniobra de escape activada por estancamiento. Con esta extensi\'{o}n el robot supera la zona y alcanza el destino.

\begin{figure}[H]
  \centering
  \begin{subfigure}[t]{0.48\linewidth}
    \centering
    \includegraphics[width=\linewidth]{nav_dijkstra_24_a_32_sin_mejora.jpg}
    \caption{Sin mejora: el robot se atasca en la zona conflictiva.}
    \label{fig:nav_24_32_sin}
  \end{subfigure}\hfill
  \begin{subfigure}[t]{0.48\linewidth}
    \centering
    \includegraphics[width=\linewidth]{nav_dijkstra_24_a_32_con_mejora.jpg}
    \caption{Con mejora: campo tangencial + escape, llega al destino.}
    \label{fig:nav_24_32_con}
  \end{subfigure}
  \caption{Comparativa en la zona conflictiva 24$\rightarrow$32: navegaci\'{o}n b\'{a}sica frente a navegaci\'{o}n con mejora anti-m\'{i}nimos locales.}
  \label{fig:nav_zona_24_32}
\end{figure}

\subsection{Reemplazo de Dijkstra por A*}
Sustituyendo Dijkstra por \texttt{aestrella(C\_eucl, H, s, t)} con la matriz de costes eucl\'{i}dea construida a partir de la adyacencia original y la heur\'{i}stica $H(i,j) = \lVert p_i - p_j\rVert_2$, A* devuelve rutas equivalentes a Dijkstra (en coste sobre $C^{\text{eucl}}$) explorando menos nodos, gracias a la informaci\'{o}n geom\'{e}trica de la heur\'{i}stica. La consistencia est\'{a} garantizada por construcci\'{o}n: tanto los pesos como la heur\'{i}stica derivan de la misma m\'{e}trica.


% =====================================================================
%  4. COMENTARIOS FINALES
% =====================================================================
\section{Comentarios finales}

La integraci\'{o}n jer\'{a}rquica de planificaci\'{o}n global y control local responde a lo que pide el enunciado: la ruta de Dijkstra (o A*) act\'{u}a como secuencia de subobjetivos, y los campos potenciales ejecutan cada tramo evitando obst\'{a}culos. Las pruebas en la zona 24--32 confirman que la versi\'{o}n b\'{a}sica falla por m\'{i}nimo local y que la mejora propuesta (campo tangencial + maniobra de escape por estancamiento) lo resuelve. El reemplazo por A* funciona correctamente al construir una matriz de costes eucl\'{i}dea coherente con la heur\'{i}stica, condici\'{o}n necesaria para que A* mantenga la garant\'{i}a de optimalidad.
